\chapter{Gérant de portefeuille et analyste de pod}\label{ch:in:portfolio-manager-and-pod-analyst}

L'une des plus grandes plateformes multigérants décrit la rémunération de ses \omterm{def:in:portfolio-manager-and-pod-analyst:roles}{gérants de portefeuille} dans le dossier qu'elle dépose auprès du régulateur américain~:
un pourcentage des profits réalisés par le gérant au cours de l'année précédente, mesurés sans tenir compte des autres gérants, avec
des pertes reportées à nouveau, un salaire qui constitue une avance sur le reversement, et aucune récupération (clawback). La plateforme alloue aussi
et réalloue le capital entre ses gérants, et peut le leur retirer. Le contrat est une option d'achat (call) sur le savoir-faire propre du gérant, assortie
d'une clause d'extinction (knock-out)~: le gérant conserve une part de la hausse, ne supporte aucune part d'une perte au-delà d'un reversement perdu, et perd son poste
si la perte devient trop profonde. Sa valeur dépend du ratio de Sharpe que le gérant possède réellement, ce que ni l'une ni l'autre des parties ne connaît
au départ. Ce chapitre décrit les métiers de \omterm{def:in:portfolio-manager-and-pod-analyst:roles}{gérant de portefeuille} et d'analyste, valorise le contrat, et se demande
ce qu'un historique de performance peut prouver.

\begin{center}\footnotesize
\begin{tabular}{@{}p{2.4cm}p{3.3cm}p{3.3cm}p{3.3cm}@{}}
\toprule
\multicolumn{4}{@{}l}{\textbf{Fiches de rôle~: le pod}}\\
\midrule
 & \omterm{def:in:portfolio-manager-and-pod-analyst:roles}{gérant de portefeuille} & \omterm{def:in:portfolio-manager-and-pod-analyst:roles}{sous-gérant de portefeuille} & \omterm{def:in:portfolio-manager-and-pod-analyst:roles}{analyste de pod}\\
\midrule
détient & le P\&L du pod & une tranche de celui-ci & rien~; il soutient le gérant\\
rémunéré par & une part du profit du pod & une part convenue avec le gérant & salaire et bonus discrétionnaire\\
rend compte à & le responsable de la stratégie de la plateforme & le \omterm{def:in:portfolio-manager-and-pod-analyst:roles}{gérant de portefeuille} & le \omterm{def:in:portfolio-manager-and-pod-analyst:roles}{gérant de portefeuille}\\
codes dans les dépôts & 11-3031, 13-2099.01 & 13-2051, 13-2099.01 & 13-2051\\
enseigné dans & livre~8, ch.~28~; livre~16, ch.~3 & livre~16, ch.~3 & livre~7, ch.~1\\
\bottomrule
\end{tabular}
\end{center}

\begin{definition}[Gérant de portefeuille, sous-gérant de portefeuille, analyste de pod]\label{def:in:portfolio-manager-and-pod-analyst:roles}
Un \emph{gérant de portefeuille}\index{gérant de portefeuille} est la personne qui décide des positions d'un portefeuille dans les
limites fixées par le propriétaire du capital, et répond de son résultat. Sur une plateforme, un \emph{sous-gérant de
portefeuille}\index{sous-gérant de portefeuille} gère une tranche du capital d'un gérant de portefeuille, avec ses propres limites et une part
de son propre profit, au sein de l'équipe du gérant. Un \emph{analyste de pod}\index{analyste de pod} travaille pour un gérant de portefeuille sur
la recherche, les données et le suivi des positions, sans capital propre.
\end{definition}

\section{Le contrat~: capital, reversement, limites et coûts}

La brochure de la plateforme est la meilleure description publique de la structure du contrat~; ses chiffres ne sont pas publiés. La
structure comporte quatre volets.
\begin{itemize}
\item \textbf{Capital.} La plateforme alloue du capital au pod (livre~8, chapitre~28), et «~allouer et
réallouer~» constitue son principal outil~: le capital augmente après de bons résultats, et diminue, voire disparaît, après de mauvais.
\item \textbf{Reversement.} Le gérant reçoit une part du profit propre du pod, au taux de reversement du livre~16, chapitre~3,
«~sans tenir compte de la performance des autres \omterm{def:in:portfolio-manager-and-pod-analyst:roles}{gérants de portefeuille}~»~; «~les pertes sont reportées à nouveau et les pertes
passées doivent être compensées avant que la rémunération liée à la performance ne devienne exigible~»~; le salaire est «~généralement
traité comme une avance~»~; et les gérants «~dont la performance est positive recevront une rémunération liée à la performance même si~» les fonds perdent de l'argent.
\item \textbf{Limites.} Limites de risque et limites de drawdown (livre~8, chapitre~28)~; l'échelle de réduction du risque (de-risking) du livre~16, chapitre~8
réduit le capital dès un premier drawdown et arrête le pod au second.
\item \textbf{Coûts.} Les coûts propres du pod (données, personnel, rémunération de l'équipe du gérant) sont imputés sur son profit avant
le reversement, et les investisseurs les supportent par les frais refacturés du livre~16, chapitre~3.
\end{itemize}
La brochure indique aussi que la plateforme a accepté de «~garantir~» un niveau de rémunération pendant un an ou plusieurs années, et de
remplacer la rémunération qu'un gérant a abandonnée en quittant un employeur précédent~: les rachats évoqués au chapitre~13.

\begin{dated}{2026-08}{Le contrat d'une plateforme avec ses gérants de portefeuille}\label{dat:in:portfolio-manager-and-pod-analyst:deal}
Partie 2A du Form ADV de Millennium Management du 25 août 2026~: rémunération «~généralement déterminée en pourcentage des profits réalisés
par ce \omterm{def:in:portfolio-manager-and-pod-analyst:roles}{gérant de portefeuille} au cours de l'année civile précédente~»~; pertes reportées à nouveau~; salaire constituant une avance sur le
reversement~; «~généralement aucun \textquotedbl carryback\textquotedbl{} ni \textquotedbl clawback\textquotedbl~»~; reversement dû même
lorsque les fonds perdent de l'argent~; garanties et rachats convenus à certaines occasions~; certaines équipes de gérants comprennent des membres
«~qui peuvent aussi être des \textquotedbl portfolio managers\textquotedbl~». Aucun taux de reversement ni seuil de drawdown n'est publié.
\end{dated}

\section{Le contrat comme option avec clause d'extinction}

\begin{method}[Le contrat du gérant de portefeuille]\label{met:in:portfolio-manager-and-pod-analyst:deal}
Soit le pod qui réalise des rendements quotidiens de ratio de Sharpe annuel $S$ et de volatilité $\sigma$ sur un capital $K$. À chaque fin d'année, le
gérant touche $\max\bigl(w,\ \rho\max(P_t-L_{t-1},0)K\bigr)$, où $P_t$ est le profit de l'année, $L_{t-1}$ la
perte reportée à nouveau, $\rho$ le taux de reversement et $w$ le salaire~; $L_t=\max(L_{t-1}-P_t,0)$. Un drawdown depuis le sommet de l'ancienneté
au-delà de $c$, réduit de moitié le capital pour le reste de l'ancienneté~; au-delà de $s$, il met fin au poste, le salaire étant versé jusqu'à
ce jour et rien après. Le gain est convexe en fonction du profit de l'année (un call) et tronqué par l'arrêt (une clause d'extinction).
\end{method}

\begin{example}[Une ancienneté de trois ans]\label{ex:in:portfolio-manager-and-pod-analyst:deal}
Prenons des conditions illustratives~: 500~millions de dollars de capital, une volatilité de 6~\% par an (30~millions de dollars), un reversement de 15~\%, un salaire de 0.5
million de dollars, un capital réduit de moitié à un drawdown de 5~\%, et le pod arrêté à 7.5~\%, sur trois ans, face à une alternative salariée
de 1.5 million de dollars par an. Sans savoir-faire ($S=0$), la rémunération espérée du gérant sur trois ans est de 3.18 millions de dollars
(médiane de 1.50 million de dollars), et le pod est arrêté dans 57.3~\% des anciennetés. Pour $S=1$, la rémunération espérée est de 10.48 millions de dollars et la
probabilité d'arrêt de 15.2~\%~; pour $S=2$, 23.36 millions de dollars et 1.9~\%. La rémunération espérée atteint l'alternative de 4.5 millions de dollars
à $S\approx0.26$, où le pod est arrêté dans environ 43~\% des anciennetés (les \cref{fig:in:portfolio-manager-and-pod-analyst:deal,fig:in:portfolio-manager-and-pod-analyst:stop}).
\end{example}

\begin{omfigure}
\begin{tikzpicture}
\begin{axis}[width=10.5cm,height=5.6cm,xmin=0,xmax=2,ymin=0,ymax=40,xlabel={\small ratio de Sharpe réel annuel du pod},
  ylabel={\small rémunération sur 3 ans, millions de \$},tick label style={font=\footnotesize},grid=major,grid style={gray!20},
  table/col sep=comma,legend style={font=\scriptsize,at={(0.5,-0.3)},anchor=north,legend cell align=left,
  legend columns=2,/tikz/every even column/.append style={column sep=8pt}}]
\addplot[name path=hi,draw=none,forget plot] table[x=sr,y=p90]{figdata/industry/22-portfolio-manager-and-pod-analyst/deal.csv};
\addplot[name path=lo,draw=none,forget plot] table[x=sr,y=p10]{figdata/industry/22-portfolio-manager-and-pod-analyst/deal.csv};
\addplot[omDef!20] fill between[of=hi and lo];
\addplot[omDef,thick,mark=*] table[x=sr,y=mean]{figdata/industry/22-portfolio-manager-and-pod-analyst/deal.csv};
\addplot[omThm,thick,dashed,mark=square*,mark options={solid}] table[x=sr,y=ce]{figdata/industry/22-portfolio-manager-and-pod-analyst/deal.csv};
\addplot[black!60,thick,domain=0:2] {4.5};
\legend{entre le 10\up{e}~centile et le 90\up{e}~centile,rémunération espérée,équivalent certain (aversion relative RRA 3),alternative salariée}
\end{axis}
\end{tikzpicture}

\omcaption{La rémunération sur trois ans du \omterm{def:in:portfolio-manager-and-pod-analyst:roles}{gérant de portefeuille} face au ratio de Sharpe réel du pod, avec les conditions illustratives de l'exemple~:
moyenne, bande comprise entre le 10\up{e}~centile et le 90\up{e}~centile, et équivalent certain pour une aversion relative au risque de 3 et 5~millions de dollars
d'autres richesses, face à une alternative salariée de 1.5 million de dollars par an. 20\,000 anciennetés par point. Données~:
\texttt{in\_pm.curve}.}\label{fig:in:portfolio-manager-and-pod-analyst:deal}
\end{omfigure}

Le contrat rémunère la chance. Sans aucun savoir-faire, sa valeur espérée est de 3.18 millions de dollars sur trois ans, soit environ les sept dixièmes
de l'alternative, car le gérant conserve une part des bonnes années et ne rend rien pendant les mauvaises, hormis les reversements
qui compensent des pertes reportées. Les investisseurs paient cette option, c'est-à-dire le coût de netting du livre~16, chapitre~3, et la plateforme la
limite par l'arrêt~: à $S=0$, plus de la moitié des anciennetés se terminent prématurément. Pour un gérant qui craint le risque, le calcul change~:
avec une aversion relative au risque de 3 et 5~millions de dollars d'autres richesses, l'équivalent certain n'atteint l'alternative que
pour $S\approx0.75$, où la probabilité d'arrêt est d'environ 23~\%. La dispersion des résultats est large à chaque niveau de savoir-faire
(la \cref{fig:in:portfolio-manager-and-pod-analyst:hist}).

\begin{omfigure}
\begin{tikzpicture}
\begin{axis}[width=10.5cm,height=4.8cm,xmin=0,xmax=2,ymin=0,ymax=60,xlabel={\small ratio de Sharpe annuel réel du pod},
  ylabel={\small anciennetés arrêtées, \%},tick label style={font=\footnotesize},grid=major,grid style={gray!20},
  table/col sep=comma]
\addplot[omThm,thick,mark=*] table[x=sr,y=stopped]{figdata/industry/22-portfolio-manager-and-pod-analyst/deal.csv};
\end{axis}
\end{tikzpicture}

\omcaption{Part des anciennetés de trois ans qui prennent fin à l'arrêt, en fonction du ratio de Sharpe réel du pod, avec l'échelle
de l'exemple (capital divisé par deux à un drawdown de 5~\%, pod arrêté à 7.5~\%). Données~: \texttt{in\_pm.curve}.}\label{fig:in:portfolio-manager-and-pod-analyst:stop}
\end{omfigure}

\begin{omfigure}
\begin{tikzpicture}
\begin{axis}[ybar interval,area legend,width=10.5cm,height=5cm,xmin=0,xmax=40,ymin=0,ymax=60,x tick label as interval=false,
  xtick={0,4,8,12,16,20,24,28,32,36,40},
  xlabel={\small rémunération sur trois ans, millions de \$ (classes de 2)},ylabel={\small part des anciennetés, \%},
  tick label style={font=\footnotesize},grid=major,grid style={gray!20},table/col sep=comma,
  legend style={font=\scriptsize,at={(0.97,0.97)},anchor=north east,legend cell align=left}]
\addplot[fill=black!25,draw=black!50] table[x=left,y=sr0]{figdata/industry/22-portfolio-manager-and-pod-analyst/hist.csv};
\addplot[fill=omDef!40,draw=omDef,fill opacity=0.6] table[x=left,y=sr1]{figdata/industry/22-portfolio-manager-and-pod-analyst/hist.csv};
\legend{aucun savoir-faire ($S=0$),$S=1$}
\end{axis}
\end{tikzpicture}

\omcaption{Distribution de la rémunération sur trois ans à deux ratios de Sharpe réels, avec les termes de l'exemple (les rémunérations supérieures à 40 millions de \$ sont rangées
dans la dernière classe). Sans savoir-faire, la plupart des anciennetés rapportent à peine plus que le salaire~; avec un ratio de Sharpe égal à un, la rémunération se répartit
du salaire à plus de 30 millions de \$. Données~: \texttt{in\_pm.run}.}\label{fig:in:portfolio-manager-and-pod-analyst:hist}
\end{omfigure}

\section{L'analyste~: de la recherche à la gestion de capital}

Un pod est une petite équipe~: un \omterm{def:in:portfolio-manager-and-pod-analyst:roles}{gérant de portefeuille}, un ou plusieurs analystes, parfois un \omterm{def:in:quant-researcher:def}{chercheur quantitatif} ou un développeur, et, dans
les pods plus importants, des \omterm{def:in:portfolio-manager-and-pod-analyst:roles}{sous-gérants de portefeuille}. La plaquette note que la supervision d'une partie des actifs d'un gérant peut être confiée à
des membres de l'équipe «~qui peuvent aussi être des \omterm{def:in:portfolio-manager-and-pod-analyst:roles}{gérants de portefeuille}~», et que certains gérants sont surtout
«~superviser et gérer d'autres membres du personnel d'investissement~». Le travail de l'analyste est la recherche du gérant~: modèles des entreprises
ou des marchés sur lesquels le pod opère, des données, le suivi des positions et du risque, et des idées.

La rémunération de l'analyste est fixée par le gérant, à partir de l'économie du pod, et elle est moins codifiée que celle du gérant. La
voie de progression consiste à gérer du capital~: d'abord une tranche du livre du gérant comme \omterm{def:in:portfolio-manager-and-pod-analyst:roles}{sous-gérant de portefeuille}, avec un historique propre, puis un
pod à soi, sur la même plateforme ou sur une autre. Cet historique est l'objet de la section suivante.

\section{Historiques de performance et mobilité}

Un gérant change de plateforme avec son historique de performance, et la plateforme qui l'embauche doit en juger. La portabilité d'un historique de performance
(livre~16, chapitre~25) est limitée par ce qu'elle peut prouver. L'arithmétique du chapitre~17 s'applique~: un ratio de Sharpe réel de 1 exige
3.85 années de résultats quotidiens pour se distinguer de zéro à 95~\%, et un historique de trois ans avec un ratio de Sharpe de 1 a une
statistique t d'environ 1.73. Un gérant arrêté la deuxième année a un historique trop court pour prouver ou réfuter quoi que ce soit, et un
gérant qui a tenu trois ans peut avoir eu de la chance~: à $S=0$, 43~\% des anciennetés de l'exemple tiennent trois ans.

La plateforme qui embauche regarde donc plus que le chiffre~: la capacité de la stratégie et son encombrement (crowding), le risque pris pour le
rendement, la part du résultat venant de quelques transactions, la possibilité d'expliquer le processus, et si l'équipe suit
le gérant. Les covenants restrictifs du chapitre~28 déterminent quand le gérant peut commencer, et le rachat du chapitre~13
fixe ce que le changement coûte au nouvel employeur.

\section{Gérants de portefeuille hors des plateformes}

Le titre de \omterm{def:in:portfolio-manager-and-pod-analyst:roles}{gérant de portefeuille} est plus ancien que les plateformes et recouvre des montages différents.
\begin{itemize}
\item \textbf{Dans un hedge fund à gérant unique} (chapitre~4), le \omterm{def:in:portfolio-manager-and-pod-analyst:roles}{gérant de portefeuille} peut être le fondateur, rémunéré par
les frais de la société et sa détention plutôt que par un taux de reversement (livre~16, chapitre~25, sur le lancement d'un fonds).
\item \textbf{Dans un fonds systématique} (chapitre~4), les portefeuilles sont gérés par des modèles~; «~\omterm{def:in:portfolio-manager-and-pod-analyst:roles}{gérant de portefeuille}~» peut désigner la personne
qui supervise un ensemble de stratégies et leur risque.
\item \textbf{Dans une société de gestion ou un fonds de pension} (chapitre~8), le \omterm{def:in:portfolio-manager-and-pod-analyst:roles}{gérant de portefeuille} gère un mandat face à un indice de référence
et est jugé par rapport à lui~; rien dans le mandat ne fonctionne comme l'arrêt de la plateforme.
\end{itemize}
Les cadres financiers de l'enquête, dont les métiers détaillés comprennent les gestionnaires de fonds d'investissement (O*NET répertorie «~Portfolio
Manager~» parmi leurs intitulés déclarés), ont perçu dans le secteur des valeurs mobilières un salaire médian de
223\,860~\$ en mai 2025, avec un 90\up{e}~centile de 383\,590~\$~; pour les analystes, 124\,370~\$ et 250\,000~\$. L'enquête ne relève
que les salaires, qui, pour un gérant de plateforme, ne représentent qu'une petite part de la rémunération.

\begin{dated}{2025-05}{Ce que gagnent les gérants de portefeuille et les analystes~: les données publiques}\label{dat:in:portfolio-manager-and-pod-analyst:pay}
Enquête professionnelle, mai 2025, secteur des valeurs mobilières~: cadres financiers (11-3031), 70\,400 salariés, entre le 10\up{e} et le 90\up{e}~centile
de 128\,260 à 383\,590~\$, médiane 223\,860~\$~; analystes financiers et d'investissement (13-2051), 89\,390, médiane 124\,370~\$, le 90\up{e}~centile
250\,000~\$. Autres fonds et véhicules d'investissement~: cadres financiers, médiane 215\,690~\$, analystes 125\,710~\$. Les
demandes de conditions de travail (LCA) pour les intitulés de \omterm{def:in:portfolio-manager-and-pod-analyst:roles}{gérant de portefeuille} chez les employeurs recensés sont trop peu nombreuses pour être publiées (14 au titre de l'exercice
fiscal 2025). Salaires uniquement.
\end{dated}

\begin{omfigure}
\begin{tikzpicture}
\begin{axis}[width=10.5cm,height=4.8cm,xmin=40,xmax=400,ymin=0.4,ymax=4.6,ytick=data,
  yticklabels={{analystes~; fonds et pools},{analystes~; valeurs mobilières},{directeurs financiers~; fonds et pools},{directeurs financiers~; valeurs mobilières}},
  yticklabel style={font=\scriptsize},xticklabel style={font=\footnotesize},xlabel={\small salaire annuel, milliers de \$},
  grid=major,grid style={gray!20},table/col sep=comma]
\addplot[draw=none,mark=none,error bars/.cd,x dir=both,x explicit,error mark=none,error bar style={omDef,line width=1pt}]
  table[x=p50,y=pos,x error plus expr=\thisrow{p90}-\thisrow{p50},x error minus expr=\thisrow{p50}-\thisrow{p10}]
  {figdata/industry/22-portfolio-manager-and-pod-analyst/survey.csv};
\addplot[draw=none,mark=none,error bars/.cd,x dir=both,x explicit,error mark=none,error bar style={omDef!40,line width=7pt}]
  table[x=p50,y=pos,x error plus expr=\thisrow{p75}-\thisrow{p50},x error minus expr=\thisrow{p50}-\thisrow{p25}]
  {figdata/industry/22-portfolio-manager-and-pod-analyst/survey.csv};
\addplot[only marks,mark=|,mark size=5pt,line width=1.5pt,omThm] table[x=p50,y=pos]{figdata/industry/22-portfolio-manager-and-pod-analyst/survey.csv};
\end{axis}
\end{tikzpicture}

\omcaption{Salaires des cadres financiers et des analystes financiers et d'investissement dans le secteur des valeurs mobilières et dans d'autres
véhicules et fonds d'investissement, enquête professionnelle de mai 2025~: du 10\up{e} au 90\up{e}~centile (trait fin), écart interquartile (trait épais),
médiane (repère). Les salaires excluent les bonus et les reversements. Données~: \texttt{data/industry/oews\_roles.csv}, au moyen de
\texttt{in\_pm.survey}.}\label{fig:in:portfolio-manager-and-pod-analyst:survey}
\end{omfigure}

\section{Tutoriel~: valoriser le contrat}\label{tut:in:portfolio-manager-and-pod-analyst:tut}

\begin{tutorial}
\textbf{Objectif.} Valoriser le contrat d'un \omterm{def:in:portfolio-manager-and-pod-analyst:roles}{gérant de portefeuille} de plateforme en fonction du savoir-faire, et trouver le savoir-faire à partir duquel il dépasse un
emploi salarié.
\textbf{État final~:} les \cref{fig:in:portfolio-manager-and-pod-analyst:deal,fig:in:portfolio-manager-and-pod-analyst:stop,fig:in:portfolio-manager-and-pod-analyst:hist}.
\begin{tutsteps}
\item \textbf{Pods.} Le \texttt{firm.podshop.pods} du livre~16 tire des rendements quotidiens pour des pods indépendants d'un ratio de Sharpe donné
et d'une volatilité donnée.
\item \textbf{Le contrat.} \texttt{firm.roles.pm\_deal} applique le reversement, les pertes reportées, l'avance sur salaire et
l'échelle sur toute l'ancienneté (voir le \cref{lst:in:portfolio-manager-and-pod-analyst:deal}).
\omcode{code/firm/roles/firm_roles.py}{347}{379}{Le contrat du gérant de portefeuille sur une ancienneté, avec une échelle de drawdown
à deux étapes.}\label{lst:in:portfolio-manager-and-pod-analyst:deal}
\item \textbf{Utilité.} La fonction \texttt{firm.payoffer.certainty\_equivalent} du chapitre~13 convertit la distribution des rémunérations en un
montant certain.
\item \textbf{Le croisement.} \texttt{in\_pm.crossing} interpole le ratio de Sharpe auquel la rémunération espérée, ou
l'équivalent certain, atteint l'alternative.
\end{tutsteps}
\end{tutorial}

Avec les paramètres de l'exemple~: la rémunération espérée atteint 4.5~millions de dollars pour $S\approx0.26$ (probabilité d'arrêt d'environ 43~\%)~;
l'équivalent certain atteint l'alternative pour $S\approx0.75$ (environ 23~\%).

\textbf{Pour aller plus loin.} Laisser la plateforme lever du capital après les bonnes années~; ajouter une garantie dès la première année~; rendre le
ratio de Sharpe incertain et laisser le gérant l'apprendre à partir des résultats~; faire supporter au pod ses coûts.

\section{Construction~: fiches du pod et contrat du gérant de portefeuille}\label{bld:in:portfolio-manager-and-pod-analyst:roles}

\begin{build}
\bfield{Rôle} Ajouter les rôles du pod au registre et modéliser le contrat du \omterm{def:in:portfolio-manager-and-pod-analyst:roles}{gérant de portefeuille}.
\bfield{Interface} \texttt{firm.roles}~: les fiches \texttt{portfolio manager}, \texttt{sub-portfolio manager},
\texttt{pod analyst}~;
\texttt{pm\_deal(sr, years, vol, capital, rate, salary, cut, stop, n, rng, days) -> dict(pay, stopped, stop\_day, pnl)}~; sur \texttt{firm.podshop.pods}.
\bfield{Règles} Reversement sur le profit propre du pod au-delà des pertes reportées~; le salaire est une avance~; aucune récupération (clawback)~; l'échelle de
réduction du risque agit sur le drawdown depuis le pic de l'ancienneté~; les termes sont fixés par l'appelant et signalés comme illustratifs.
\bfield{Tests d'acceptation} \texttt{code/firm/roles/tests/}~: un pod qui ne peut pas perdre verse exactement le taux multiplié par son profit~;
un arrêt dès le premier jour verse un jour de salaire~; davantage de compétence signifie moins d'arrêts et davantage de rémunération.
\bfield{Pour aller plus loin} Capital qui croît avec les résultats~; un gérant qui cesse de trader après une année maximisant le reversement~; le
côté de la plateforme~: reversement espéré comparé au rendement net des investisseurs.
\end{build}

\begin{omsources}
\item Millennium Management LLC, Form ADV Part 2A brochure, 25 août 2026 (SEC IAPD).
\item Bureau of Labor Statistics, enquête sur les professions, mai 2025~; tableaux LCA du chapitre~14.
\item livre~16, chapitres~3, 8 et~25~; livre~8, chapitre~28.
\end{omsources}

\section{Exercices}

\begin{exercise}[$\star$]\label{exo:in:portfolio-manager-and-pod-analyst:1}
Un pod a réalisé 40~millions de dollars l'an dernier, après une perte de 10~millions de dollars l'année précédente. Avec un taux de reversement de 15~\% et un salaire de 0.5~million
de dollars, combien le gérant a-t-il perçu pour l'année dernière~?
\end{exercise}

\begin{exercise}[$\star$]\label{exo:in:portfolio-manager-and-pod-analyst:2}
D'après l'exemple, de combien la rémunération espérée sur trois ans augmente-t-elle entre $S=0$ et $S=1$, et de combien
baisse la probabilité d'arrêt~?
\end{exercise}

\begin{exercise}[$\star$]\label{exo:in:portfolio-manager-and-pod-analyst:3}
Quelle t-statistique présente un historique quotidien de trois ans dont le ratio de Sharpe vaut 2~?
\end{exercise}

\begin{exercise}[$\star\star$]\label{exo:in:portfolio-manager-and-pod-analyst:4}
Pourquoi le contrat vaut-il 3.18~millions de dollars en espérance pour un gérant sans compétence, et qui en supporte le coût~?
\end{exercise}

\begin{exercise}[$\star\star$]\label{exo:in:portfolio-manager-and-pod-analyst:5}
Pourquoi l'équivalent certain croise-t-il l'alternative pour un ratio de Sharpe plus élevé que la rémunération espérée~?
\end{exercise}

\begin{exercise}[$\star\star$]\label{exo:in:portfolio-manager-and-pod-analyst:6}
Un analyste se voit proposer un sous-portefeuille de 50~millions de dollars avec une part de 10~\% de son profit, ou un salaire de 300\,000~dollars avec un
bonus discrétionnaire. Que doit-elle savoir pour choisir~?
\end{exercise}

\begin{exercise}[$\star\star\star$]\label{exo:in:portfolio-manager-and-pod-analyst:7}
\emph{Programmation.} Relancer l'exemple avec l'arrêt déplacé de 7.5~\% à 10~\%. Que deviennent la rémunération espérée et la
probabilité d'arrêt en $S=0$ et $S=1$~?
\end{exercise}

\begin{exercise}[$\star\star\star$]\label{exo:in:portfolio-manager-and-pod-analyst:8}
\emph{Trouver la faille.} «~Elle a dirigé un pod pendant trois ans avec un ratio de Sharpe de 1.2 et n'a jamais été arrêtée~: elle est manifestement
compétente, et il faudrait lui donner deux fois plus de capital.~»
\end{exercise}

\section{Problème~: Le contrat}

\begin{problem}[{Problème du week-end --- le contrat}]\label{pb:in:portfolio-manager-and-pod-analyst:1}
Un analyste senior dans une société de gestion, payé 1.5~million de dollars par an, se voit proposer un pod sur une plateforme aux conditions de l'exemple.
Elle croit que son véritable ratio de Sharpe se situe entre 0.5 et 1.

\medskip\noindent\textbf{Partie I --- Les rôles et le contrat.}
\begin{enumerate}
\item Définir le \omterm{def:in:portfolio-manager-and-pod-analyst:roles}{gérant de portefeuille}, le \omterm{def:in:portfolio-manager-and-pod-analyst:roles}{sous-gérant de portefeuille} et l'\omterm{def:in:portfolio-manager-and-pod-analyst:roles}{analyste de pod}.
\item Décrire la structure de reversement que la brochure de la plateforme énonce.
\item Que fait la plateforme du capital, et pourquoi est-ce son principal outil~?
\item Qu'est-ce que l'échelle de réduction du risque (de-risking), et quel livre la définit~?
\item Quels frais sont prélevés avant le reversement~?
\end{enumerate}

\medskip\noindent\textbf{Partie II --- Le modèle.}
\begin{enumerate}[resume]
\item Énoncer la fonction de gain du contrat, avec pertes reportées et avance sur salaire.
\item Pourquoi s'agit-il d'un call, et pourquoi d'un knock-out~?
\item Énoncer les paramètres de l'exemple.
\item Donner la rémunération espérée et la probabilité d'arrêt en $S=0$, 1 et 2.
\item Qu'est-ce que l'équivalent certain, et pour quelle aversion au risque et quelle richesse est-il calculé~?
\end{enumerate}

\medskip\noindent\textbf{Partie III --- Historiques et alternatives.}
\begin{enumerate}[resume]
\item Pendant combien de temps un ratio de Sharpe de 1 doit-il se maintenir pour se distinguer de zéro~?
\item Que lit la plateforme de recrutement au-delà du chiffre~?
\item En quoi le contrat d'un \omterm{def:in:portfolio-manager-and-pod-analyst:roles}{gérant de portefeuille} dans une société de gestion diffère-t-il~?
\item Que montrent les salaires de l'enquête, et que laissent-ils de côté~?
\item Comment un analyste évolue-t-il vers la gestion d'un capital~?
\end{enumerate}

\medskip\noindent\textbf{Partie IV --- Le verdict.}
\begin{enumerate}[resume]
\item Énoncer le \emph{résultat nommé}~: le ratio de Sharpe véritable à partir duquel la rémunération espérée sur trois ans dépasse l'alternative salariée
et la probabilité d'arrêt à ce niveau.
\item Que devient la réponse avec son aversion au risque~?
\item Sur l'intervalle de ses croyances, quelles sont la rémunération espérée et la probabilité d'arrêt~?
\item Que doit-elle négocier~?
\item En deux phrases, lui donner un conseil.
\end{enumerate}
\end{problem}

\section{Questions d'entretien}

\begin{interviewq}[$\star$ \iqroles{trader, chercheur}]\label{iq:in:portfolio-manager-and-pod-analyst:1}
Votre pod est en baisse de 4~\% par rapport à son pic, avec un arrêt à 7.5~\%. Que faites-vous~?
\end{interviewq}

\begin{interviewq}[$\star$ \iqroles{chercheur}]\label{iq:in:portfolio-manager-and-pod-analyst:2}
Comment annonceriez-vous à un \omterm{def:in:portfolio-manager-and-pod-analyst:roles}{gérant de portefeuille} que votre meilleure idée ne fonctionne plus~?
\end{interviewq}

\begin{interviewq}[$\star\star$ \iqroles{trader}]\label{iq:in:portfolio-manager-and-pod-analyst:3}
Pourquoi un gérant avec des pertes reportées pourrait-il prendre davantage de risque, et comment une plateforme pourrait-elle l'en empêcher~?
\end{interviewq}

\begin{interviewq}[$\star\star$ \iqroles{chercheur, trader}]\label{iq:in:portfolio-manager-and-pod-analyst:4}
Quelle part du rendement de votre pod est due au bêta de marché, et comment la mesureriez-vous~?
\end{interviewq}

\begin{interviewq}[$\star\star$ \iqroles{trader}]\label{iq:in:portfolio-manager-and-pod-analyst:5}
Que voudriez-vous savoir de l'historique sur trois ans d'un gérant candidat avant de le recruter~?
\end{interviewq}

\begin{interviewq}[$\star\star\star$ \iqroles{chercheur, trader}]\label{iq:in:portfolio-manager-and-pod-analyst:6}
Évaluez l'option du gérant sur un an~: le profit du pod suit une loi normale de moyenne nulle et d'écart-type 30~millions de dollars, le
reversement de 15~\%. Quel est le reversement espéré~?
\end{interviewq}
